\documentclass[11pt,a4paper]{article}
\usepackage[margin=0.85in,top=1in,bottom=1in]{geometry}
\usepackage{amsmath,amssymb,amsfonts}
\usepackage{graphicx}
\usepackage{float}
\usepackage{enumitem}
\usepackage{microtype}
\usepackage{caption}
\usepackage[hidelinks]{hyperref}
\usepackage[most,breakable]{tcolorbox}
\tcbuselibrary{breakable,skins}
\usepackage[utf8]{inputenc}
\usepackage[T1]{fontenc}
\usepackage{lmodern}
\captionsetup{font=small,labelfont=bf,skip=4pt}
\setlist{leftmargin=1.5em,topsep=2pt}

%% Breakable card — prevents content cutoff across pages
\newtcolorbox{solcard}[3]{
  breakable, enhanced,
  colback=white, colframe=gray!50, arc=2pt,
  title={\textbf{#1}\hfill\textsf{\small #3\,pts}},
  fonttitle=\normalsize\sffamily\bfseries,
  before upper={\small\textit{#2}\par\smallskip\hrule height 0.4pt\smallskip},
  top=4pt, bottom=6pt, left=7pt, right=7pt,
  parbox=false,
}

\title{\textbf{Faraday effect in electron gas --- Solution}}
\author{\normalsize Y26 (2026) — English translation}
\date{}
\begin{document}\maketitle\vspace{-0.5em}
\begin{solcard}{A1}{Show that the dependence of conductivity on frequency in the model under consideration has the form $$ \sigma = \frac{D}{-i \omega + \frac{1}{\tau}}. $$Express coefficient $D$ through $n$, $q$, $m$.}{0.30}

The equation of electron motion is $$ m \dot{\vec{v}} = - \frac{m}{\tau} \vec{v}  - q \vec{E}. $$Substitute solution in form exponential, obtain $$ m\vec{v}_0\left( -i \omega + \frac{1}{\tau}\right) = - q \vec{E}_0, $$from which $$ \vec{v_0} = \frac{-q\vec{E}_0/m}{-i \omega + 1/\tau}. $$Then density current $$ \vec{j} = -q n \vec{v} = \frac{n q^2/m}{-i \omega + 1/\tau} \vec{E}, $$and conductivity $$ \sigma = \frac{n q^2/m}{-i \omega + 1/\tau}. $$

\par\smallskip\noindent\textbf{Answer:} $$ D = \frac{n q^2}{m} $$
\end{solcard}

\begin{solcard}{A2}{Let now the system be placed in a magnetic field $B$ directed along the axis $z$. The complex amplitude of the electric field $\vec{E}_0$ lies in the $xy$ plane, its projections $E_{0x}$, $E_{0y}$. Write down the equations to determine the components of the steady-state speed of electron motion $v_{0x}$, $v_{0y}$.}{0.50}

The equations of motion taking into account the magnetic field have the form $$ m \dot{\vec{v}} = - \frac{m}{\tau} \vec{v}  - q( \vec{E} + \vec{v} \times \vec{B}). $$Substituting the solution in exponential form and writing the vector product through the components, we obtain

\par\smallskip\noindent\textbf{Answer:} \begin{align*} m v_{0x}\left( -i \omega + \frac{1}{\tau}\right) = - q \left(E_{0x}+ v_{0y}B \right), \\ m v_{0y}\left( -i \omega + \frac{1}{\tau}\right) = - q \left(E_{0y} - v_{0x} B\right) \\ \end{align*}
\end{solcard}

\begin{solcard}{A3}{Show that conductivities in a magnetic field satisfy the relations $\sigma_{xx} = \sigma_{yy}$, $\sigma_{xy} = - \sigma_{yx}$. Express $\sigma_{xx}$ and $\sigma_{xy}$ in terms of $\omega$, $D$, $\tau$ and the frequency of electron motion in the magnetic field $\omega_B = qB/m$.}{1.20}

Let's move the terms with speed to the left side, divide the equations by mass and introduce frequency $\omega_B$ instead of the magnetic field: \begin{align*} \left( -i \omega + \frac{1}{\tau}\right) v_{0x} +\omega_B v_{0y} = - \frac{q}{m} E_{0x}, \\ - \omega_B v_{0x} +\left( -i \omega + \frac{1}{\tau}\right) v_{0y} = - \frac{q}{m} E_{0y} . \\ \end{align*} Multiply the first equation by $(-i\omega + 1/\tau)$, the second by $-\omega_B$ and add, we get $$ \left( (-i\omega + 1/\tau)^2 + \omega_B^2\right) v_{0x} = - \frac{q}{m}\left((-i\omega + 1/\tau) E_{0x} - \omega_B E_{0y} \right). $$Similarly first equation multiply by $\omega_B$, the second by $(-i\omega + 1/\tau)$: $$ \left( (-i\omega + 1/\tau)^2 + \omega_B^2\right) v_{0y} = - \frac{q}{m}\left(\omega_B E_{0x} + (-i\omega + 1/\tau) E_{0y}  \right). $$ Use relation density current and velocity $\vec{j} = -q n \vec{v}$, find (exponential factor identical, therefore for current hold such same relation, as and for amplitude) $$ j_x = \frac{nq^2}{m} \frac{(-i\omega + 1/\tau) E_{x} - \omega_B E_{y} }{ (-i\omega + 1/\tau)^2 + \omega_B^2}, \quad j_y = \frac{nq^2}{m} \frac{\omega_B E_{x} + (-i\omega + 1/\tau) E_{y}  }{ (-i\omega + 1/\tau)^2 + \omega_B^2}. $$Note, that for zero magnetic field $\omega_B = 0$ and be obtained relation from previous part. Since this expression coincides with the one given in the condition, from it one can obtain expression for conductivity, equality $\sigma_{xx} = \sigma_{yy}$ and $\sigma_{xy} = - \sigma_{yx}$ follow directly from these expressions.

\par\smallskip\noindent\textbf{Answer:} \begin{align*} \sigma_{xx} = \sigma_{yy} = \frac{D (-i \omega + 1/\tau)}{(-i\omega + 1/\tau)^2 + \omega_B^2},\\ \sigma_{xy} = - \sigma_{yx} = -\frac{D \omega_B}{(-i\omega + 1/\tau)^2 + \omega_B^2} \end{align*}
\end{solcard}

\begin{solcard}{B1}{Find the ratio of complex field amplitudes $E_{0y}/E_{0x}$ for right and left circular polarization.}{0.30}

A field with right circular polarization rotates in the $xy$ plane counterclockwise, in real form the dependence has the form $$ E_{x} = A \cos \omega t, \quad E_{y} = A \sin \omega t. $$Going over to complex notation, we obtain $$ E_{x} = \mathrm{Re}(A e^{-i \omega t}), \quad E_{y} = \mathrm{Re}(iA e^{-i \omega t}). $$Therefore for right polarization $E_{0y}/E_{0x} = i$. For left polarization the opposite direction of rotation can be achieved by changing the sign of $E_y$.

\par\smallskip\noindent\textbf{Answer:} Right polarization: $E_{0y}/E_{0x} = i$, left polarization: $E_{0y}/E_{0x} =- i$
\end{solcard}

\begin{solcard}{B2}{Show that the relations $\vec{j} = \sigma_{\pm} \vec{E}$ hold for circular polarizations and find the corresponding values of conductivity $\sigma_{\pm}$. Express your answer in terms of conductivities $\sigma_{xx}$ and $\sigma_{xy}$.}{0.50}

First, let’s substitute a field with right polarization into the formula for the current density, that is, $E_{0x} = A$, $E_{0y }= iA$: $$ j_{0x} = \sigma_{xx} A + i \sigma_{xy} A, \quad j_{0y} = - \sigma_{xy} A + i \sigma_{xx} A. $$Here in second equality use relation between conductivity. Hence we get $j_{0y} = i j_{0x}$, and means $$ \vec{j} = (\sigma_{xx} + i \sigma_{xy}) \vec{E}, \quad \sigma_{+} =\sigma_{xx} + i \sigma_{xy}. $$An analogous calculation for left polarization differ only sign before $\sigma_{xy}$.

\par\smallskip\noindent\textbf{Answer:} $$ \sigma_{+} =\sigma_{xx} + i \sigma_{xy}, \quad \sigma_{-} =\sigma_{xx} - i \sigma_{xy} $$
\end{solcard}

\begin{solcard}{B3}{For a medium with the relation $\vec{j} = \sigma(\omega) \vec{E}$, determine the dielectric constant $\varepsilon(\omega)$ and the corresponding refractive index $n$, express the answer in terms of $\sigma$, $\omega$ and fundamental constants.}{0.40}

The electric induction vector is expressed through the electric field and the polarization vector by the relations $$ \vec{D} = \varepsilon_0 \varepsilon \vec{E}  = \varepsilon_0 \vec{E} + \vec{P}. $$Hence $$ \vec{P} = \varepsilon_0 (\varepsilon - 1) \vec{E}. $$Express density current through the derivative of the polarization vector and compare with the expression through the conductivity $$ \vec{j} = - i \omega  \varepsilon_0 (\varepsilon - 1) \vec{E} = \sigma \vec{E}. $$From here we get

\par\smallskip\noindent\textbf{Answer:} $$ \varepsilon = 1 + \frac{i \sigma}{\varepsilon_0 \omega}, \quad n = \sqrt{\varepsilon} = \sqrt{1 + \frac{i \sigma}{\varepsilon_0 \omega}} $$
\end{solcard}

\begin{solcard}{B4}{Assuming that the refractive indices are close to $1$, find the difference in refractive indices for waves with right and left polarization $\Delta n = n_{+} - n_{-}$. Express your answer in terms of $\omega$, $\sigma_{xx}$, $\sigma_{xy}$ and the fundamental constants.}{0.50}

Since the refractive index is close to 1, the expression from the previous paragraph can be expanded to the first order, for circular polarizations we get $$ n_{\pm} = 1 + \frac{i \sigma_{\pm}}{2 \varepsilon_0 \omega} = 1 + \frac{i (\sigma_{xx} \pm i \sigma_{xy})}{2 \varepsilon_0 \omega}. $$Compute difference, find $$\Delta n  = \frac{i (\sigma_{xx} + i \sigma_{xy})}{2 \varepsilon_0 \omega} - \frac{i (\sigma_{xx} - i \sigma_{xy})}{2 \varepsilon_0 \omega} = - \frac{\sigma_{xy}}{\varepsilon_0 \omega}.$$Note, that conductivity $\sigma_{xx}$ cancels to first order.

\par\smallskip\noindent\textbf{Answer:} $$ \Delta n = - \frac{\sigma_{xy}}{\varepsilon_0 \omega}. $$
\end{solcard}

\begin{solcard}{B5}{Let a plane monochromatic wave with linear polarization fall on a layer of medium of thickness $h$. The boundaries of the medium are perpendicular to the $z$ axis, the wave propagates along the $z$ axis. Show that after passing through the medium, the direction of polarization rotates through an angle $\theta_F$. Express this angle in terms of $h$, $\sigma_{xy}$ and fundamental constants.}{0.80}

Without loss of generality, we can assume that initially the wave is polarized along the $x$ axis, the wave amplitude is $A$. Let's imagine it in the form of a superposition of two waves with circular polarizations, we will write the vectors of complex amplitudes in the form of columns: $$ \vec{E}_{0} = \begin{pmatrix} A \\ 0\end{pmatrix} = \frac{A}{2} \begin{pmatrix} 1 \\ i \end{pmatrix} + \frac{A}{2} \begin{pmatrix} 1 \\ -i \end{pmatrix}. $$After passing through the medium, different circular polarizations acquire different phases, amplitude field after passage through the medium $$ \vec{E}_1 = \frac{A}{2} \begin{pmatrix} 1 \\ i \end{pmatrix} e^{i k n_+ h}+ \frac{A}{2} \begin{pmatrix} 1 \\ -i \end{pmatrix} e^{i k n_- h}. $$Here $k$ – wave number in vacuum. Factoring out the common phase $\Delta \varphi_0 = kh (n_+ + n_-)/2$, obtain $$ \vec{E}_1 = e^{i \Delta \varphi_0} \frac{A}{2} \begin{pmatrix} e^{i k \Delta n h/2} +  e^{-i k \Delta n h/2}\\ i (e^{i k \Delta n h/2} -  e^{-i k \Delta n h/2})\end{pmatrix} = A e^{i \Delta \varphi_0} \begin{pmatrix} \cos \frac{k h \Delta n}{2} \\ - \sin\frac{k h \Delta n}{2} \end{pmatrix}. $$Obtain expression for linear polarization, rotate on angle $$ \theta_F = - \frac{k \Delta n h}{2}. $$The minus sign means that the rotation occurs clockwise. Substitute expression for difference exponent/index refraction and use relation for frequency and wave number in vacuum $\omega = c k$, obtain $$ \theta_F = \frac{k h \sigma_{xy}}{2 \varepsilon_0 \omega} = \frac{\sigma_{xy} h}{2 \varepsilon_0 c}. $$


The same result for the rotation angle can be obtained using vector diagrams. A wave with linear polarization can be represented as a superposition of two waves with circular polarizations and equal amplitudes. These waves can be represented as a pair of vectors in the $xy$ plane rotating counterclockwise and clockwise, respectively. In this case, oscillations in a wave with linear polarization occur along the bisector of the angle between two vectors of circular polarizations, initially this is the $x$ axis. After passing through the medium, an additional phase difference $\Delta \varphi$ arises between waves with circular polarizations. This means that for the same position of the wave vector with left polarization, the wave vector with right polarization will rotate by $\Delta \varphi = - k h \Delta n$. The sign $-$ is due to the fact that the oscillation has the form $$\operatorname{Re} e^{-i \omega t + i kh\Delta n} = \cos (\omega t - \Delta \varphi).$$ Resultant oscillation still still represent itself sum two wave with circular polarization and equal amplitude, but the bisector between their vectors rotates by the angle $\theta_F = \Delta \varphi/2$, which coincides with the answer found earlier.

\par\smallskip\noindent\textbf{Answer:} $$ \theta_F = \frac{\sigma_{xy} h}{2 \varepsilon_0 c} $$
\begin{figure}[H]
  \centering
  \includegraphics[width=0.72\textwidth]{T5_S_EN_assets/fig1.jpg}
\end{figure}
\end{solcard}

\begin{solcard}{B6}{Using the result of step $\textbf{A3}$, obtain a real expression for the angle of rotation of the plane of polarization $\theta_F$. Leave the real part of the expression obtained above. Express your answer in terms of $D$, $h$, $\tau$, $\omega_B$, $\omega$ and fundamental constants.}{0.50}

Let's substitute the expression for conductivity: $$ \theta_F = - \frac{D \omega_B h}{2 \varepsilon_0 c}\frac{1}{(-i\omega + 1/\tau)^2 + \omega_B^2}. $$Over the rotation of the plane of polarization corresponds to the real part of this expression, so to extract it, we expand the bracket in the denominator and multiply numerator and denominator by the conjugate expression: $$ \frac{1}{(-i\omega + 1/\tau)^2 + \omega_B^2} = \frac{1}{\omega_B^2 + 1/\tau^2- \omega^2 - 2 i \omega/\tau} = \frac{\omega_B^2 + 1/\tau^2- \omega^2 + 2 i \omega/\tau}{(\omega_B^2 + 1/\tau^2- \omega^2)^2 +4 \omega^2/\tau^2}. $$Taking the real part, we get

\par\smallskip\noindent\textbf{Answer:} $$ \theta_F = - \frac{D \omega_B h}{2 \varepsilon_0 c}\frac{\omega_B^2 + 1/\tau^2- \omega^2 }{(\omega_B^2 + 1/\tau^2- \omega^2)^2 +4 \omega^2/\tau^2}. $$
\end{solcard}

\begin{solcard}{B7}{In general, the refractive indices for both circular polarizations will contain different imaginary parts, which will generally lead to an elliptical polarization. Let the parameters of the system be such that $kn_+h = 0.73 +0.18i$, $k n_- h = 1.17 + 0.35i$, where unimportant terms that are multiples of $2\pi$ are discarded. Here $k$ is the wave number in vacuum. The polarization of the incident wave is directed along the $x$ axis. Find the ratio of the semi-axes $a/b$ and the angle $\psi$ between the semi-major axis and the $x$ axis for an ellipse of finite polarization.}{1.50}

As in the previous part, the difference in the real parts of the products $kh (n_+ - n_-)$ leads to a change in the direction of polarization, with the rotation angle $$ \psi_1 = - \frac{k h}{2}\mathrm{Re} \Delta n = 0.22 ~\text{rad} \approx 12.6^\circ. $$Choose direction new axis $x'$ along rotate direction polarization, and $y'$ perpendicularly it. The difference of imaginary components leads to a change in amplitude of the circular polarization, therefore in rotate axes the complex amplitude has the form (ignoring an inessential common phase) $$ \vec{E}_1 = \frac{A}{2} \begin{pmatrix} 1\\ i \end{pmatrix} e^{- \gamma_1} +  \frac{A}{2} \begin{pmatrix} 1\\ -i \end{pmatrix} e^{-\gamma_2} = \frac{A}{2} \begin{pmatrix} e^{- \gamma_1} + e^{-\gamma_2} \\ i(e^{-\gamma_1} - e^{-\gamma_2})\end{pmatrix}. $$Here exponent/index exponential $\gamma$ be determined the imaginary part expression from the condition: $\gamma_1 = k h \mathrm{Im}\, n_+ = 0.18$, $\gamma_2 = k h \mathrm{Im}\, n_- = 0.35$.  Shear phase between oscillation by axis $x'$ and $y'$ after accounting for damping still still $\pi/2$, therefore this axis coincides with the principal axes of the polarization ellipse. Means ratio axis $$ \frac{a}{b} = \frac{e^{-\gamma_1} + e^{- \gamma_2}}{e^{-\gamma_1} - e^{- \gamma_2}} \approx 11.8, $$and larger axis rotate on angle $\psi = \psi_1$ relative to the original axis.

\par\smallskip\noindent\textbf{Answer:} $$ \psi \approx 0.22 ~\text{rad} \approx 12.6^\circ, \quad \frac{a}{b} \approx 11.8, \quad \frac{b}{a} \approx 0.0845 $$
\end{solcard}

\begin{solcard}{C1}{Find the frequency $\omega_B$ of motion of an electron with energy $\varepsilon_F$ in a magnetic field. Express your answer using $v_F$, $q$, $B$, $\varepsilon_F$.}{0.50}

Since the standard expression for the Lorentz force is applicable, the magnetic force acting on the electron will is equal to in magnitude to $qv_F B$ and directed perpendicular to the velocity. The modulus of the electron's momentum remains constant due to the law of conservation of energy, therefore the movement occurs in a circle and the law of change in momentum has the form $$ \omega_B p= q v_F B. $$Substitute expression for magnitude/modulus momentum $p = \varepsilon_F/v_F$, we find

\par\smallskip\noindent\textbf{Answer:} $$ \omega_B = \frac{q v_F^2 B}{\varepsilon_F} $$
\end{solcard}

\begin{solcard}{C2}{After analyzing the given graph, determine the parameters $D_2$, $\tau$ and $\omega_B$ for the graphene sample under study. Magnetic field $B = 7~\text{T}$. When analyzing the graph, we can assume that $\omega_B \tau \gg 1$.}{1.50}

We will analyze the graph using the formula from $\mathbf{B6}$ with the replacements indicated in the condition $$ \theta_F =  C\frac{1 +\tau^2( \omega_B^2 - \omega^2) }{(1 +\tau^2( \omega_B^2 - \omega^2))^2 +4 \omega^2\tau^2}, \quad C = \frac{D_2 \omega_B  \tau^2}{2 \varepsilon_0 c}. $$The overall sign is inessential, since the direction of rotation is not specified. in force condition $\omega_B^2 \tau \gg 1$ in numerator and in first bracket in denominator one can neglect 1: $$ \theta_F =  C\frac{ \omega_B^2 - \omega^2}{\tau^2(\omega^2 - \omega_B^2)^2 +4 \omega^2}. $$ Angle rotation bebecomes 0 for $\omega^2 = \omega_B^2$. From the graph it can be seen, that in region negative $\theta_F$ is maximum magnitude/modulus angle rotation. Find it position analytically. For this introduce auxiliary variable $y = \omega_B^2 - \omega^2$. Angle rotation be expressed through it as $$ \theta_F = C \frac{y}{y^2 \tau^2 + 4\omega_B^2 - 4 y} = \frac{C}{y \tau^2 + \frac{4 \omega_B^2}{y}- 4}. $$Denominator minimum for $y = \pm 2 \omega_B/\tau$, moreover for the maximum modulus on the graph, $\omega > \omega_B$ and we should choose the negative root. Correspond value angle $$ \theta_1 =- \frac{C}{ 4 (\omega_B \tau + 1)} $$From the graph, 0 and the maximum modulus correspond to frequency $$ f_0 \approx 4.7 ~\text{THz}, \quad f_{\text{max}} \approx 6.5~\text{THz}. $$Hence immediately find $\omega_B = 2\pi f_0 \approx 29.5 \cdot10^{12} ~\text{s}^{-1}$. From position maximum $$ (2\pi f_{\text{max}})^2 = \omega_B^2 + \frac{2 \omega_B}{\tau}, $$hence find $$ \frac{1}{\tau} = 13.5\cdot 10^{12} ~\text{c}^{-1}, \quad \tau = 7.4 \cdot 10^{-14} ~\text{s}. $$The value $\omega_B \tau \approx 2$ is not very large, however for calculation position the maximum the answer contains the square of this expression, which gives an error order $20\%$. From value in maximum $\theta_1 \approx -0.055$ find $$C \approx 0.70,$$and we obtain the Drude weight $$ D_2 = \frac{2 \varepsilon_0 c}{\omega_B \tau^2} C \approx 2.3\cdot 10^{10} ~\text{C} ^2\cdot \text{kg}^{-1} \cdot \text{m}^{-2}. $$

\end{solcard}

\begin{solcard}{C3}{Using the results from the previous paragraph, determine the values of $\varepsilon_F$ (in electronvolts) and $v_F$.}{0.50}

Using the value $D_2$ we find $$ \varepsilon_F = \frac{2\hbar^2}{q^2} D_2 \approx 0.125 ~\text{eV}. $$From the formula for frequency motion in magnetic field $$ v_F = \sqrt{\frac{\varepsilon_F \omega_B}{q B}} \approx 0.7 \cdot 10^6~\text{m}/\text{s}. $$

\end{solcard}

\begin{solcard}{D1}{Let natural light fall on the system on the left with intensity $I_0$. Find the intensity $I_1$ of radiation at the exit from the system. Find also the intensity $I_2$ at the output if the same natural light falls on the system from the right. At what values of $B_z$ will the intensity ratio $I_2/I_1$ be minimal?}{1.00}

After passing through the first polaroid, the intensity decreases by a factor of 2, and the polarization is directed along the $x$ axis. When moving along the $z$ axis, the polarization rotates through an angle $\varphi$ to the $y$ axis, and the intensity of the transmitted light will is equal to to $$ I_1 = \frac{I_0}{2} \cos^2 (\psi - \varphi) = \frac{I_0}{2} \cos^2 (\psi - V h B). $$For motion in inverse side direction rotation polarization change, therefore the expression for the intensity $$ I_2 = \frac{I_0}{2} \cos^2 (\psi + \varphi) = \frac{I_0}{2} \cos^2 (\psi + V h B). $$The ratio $I_2/I_1$ is minimal when the transmitted light is absent and $I_2 = 0$. This be achieved for $$ \psi + Vh B = \frac{\pi }{2} + \pi m, $$

\par\smallskip\noindent\textbf{Answer:} $$ B = \frac{\psi - \pi/2 - \pi m}{V h}. $$
\end{solcard}

\end{document}